% !TEX root = ../main.tex
\section{Evaluation}
\label{sec:eval}

We attribute \sys's speedup (\S\ref{sec:eval-e2e}), examine how capacities and aggregation interact
(\S\ref{sec:eval-abl}--\S\ref{sec:eval-pad}), and test blocks, multiple chunks, generalization, accuracy and a fused kernel
(\S\ref{sec:eval-blocks}--\S\ref{sec:eval-kernel}).

\subsection{Experimental Setup}
\label{sec:setup}

\textbf{Hardware and model.} One Cloud AI 100 Ultra card, four SoCs with 16 cores and 32\,GB each, platform SDK 1.21.6;
Qwen3-30B-A3B, all 48 layers, batch 1, MXFP6 matrix weights and FP16 activations as in the production deployment, and the
KV cache on device.

\textbf{Programs.} The ladder of Table~\ref{tab:ladder} adds one step at a time: (1) the production compilation of the
exported graph; (2) expert-parallel compilation with head-parallel attention, the export's combine and capacity $T$, our
\emph{naive \naive} baseline; (2$'$) attention's sum over its head groups written elementwise, kept by every later
program; (3a) tree prefix scans for the routing; (3b) two repairs of the dense combine; (4) token-owned aggregation and
(5) its elementwise cross-SoC sum (\S\ref{sec:red}); (6) the run-time rank sort with two tiers at $T$ and $T/8$; (7) \sys's
three tiers. Controlled variants change one choice of these programs.

\textbf{Inputs and protocol.} At $T=128$ the first chunk of a GSM8K chat prompt, at 256 and 512 held-out GSM8K chunks
outside the calibration range; latency depends on a chunk only through its routing, and on one layer it varied by about
1\% across 200 prompts of four workloads. A pass is the median of three rounds of 20 inferences after warm-up. Sessions
time two to six programs in two or three passes of mirrored order, always with program 5 as an anchor, and one
program$\,+\,$session fit of log latency over all passes (66 to 75 per chunk length) gives each program's latency; the
largest residual is 0.50\%, 0.88\% and 1.7\% at $T=128$, 256 and 512. We compare the last-position logits with an FP32 CPU run and read back every lane's count.

\subsection{End-to-End Results}
\label{sec:eval-e2e}

\begin{table}[t]
\caption{Prefill latency of one chunk (ms) along the ladder, the change against the step above, and the contribution of
each group of steps.}
\label{tab:ladder}
\vskip 0.1in
\begin{center}
\begin{footnotesize}
\setlength{\tabcolsep}{1.2pt}
\begin{tabular}{@{}lr@{\,}lr@{\,}lr@{\,}l@{}}
\toprule
Program & \multicolumn{2}{c}{$T=128$} & \multicolumn{2}{c}{$T=256$} & \multicolumn{2}{c}{$T=512$} \\
\midrule
1 Production              & 499.6 &  & 516.5 &  & 1053.5 &  \\
2 + expert parallel       & 266.9 & {\scriptsize $-$46.6\%} & 569.3 & {\scriptsize $+$10.2\%} & 1157.6 & {\scriptsize $+$9.9\%} \\
2$'$ + head-group sum     & 250.0 & {\scriptsize $-$6.3\%} & 502.4 & {\scriptsize $-$11.8\%} & 1027.4 & {\scriptsize $-$11.2\%} \\
3a + tree scans           & 212.4 & {\scriptsize $-$15.1\%} & 431.0 & {\scriptsize $-$14.2\%} & 877.8 & {\scriptsize $-$14.6\%} \\
3b + dense repairs        & 141.9 & {\scriptsize $-$33.2\%} & 289.9 & {\scriptsize $-$32.7\%} & 570.9 & {\scriptsize $-$35.0\%} \\
4 + token-owned           & 115.8 & {\scriptsize $-$18.5\%} & 189.9 & {\scriptsize $-$34.5\%} & 350.9 & {\scriptsize $-$38.5\%} \\
5 + elementwise sum       & 113.2 & {\scriptsize $-$2.2\%} & 184.8 & {\scriptsize $-$2.7\%} & 349.2 & {\scriptsize $-$0.5\%} \\
6 + rank sort, 2 tiers    & 106.3 & {\scriptsize $-$6.1\%} & 138.0 & {\scriptsize $-$25.3\%} & 269.6 & {\scriptsize $-$22.8\%} \\
7 + 3 tiers (\sys)        & \textbf{105.3} & {\scriptsize $-$0.9\%} & \textbf{141.5} & {\scriptsize $+$2.5\%} & \textbf{225.3} & {\scriptsize $-$16.4\%} \\
\midrule
Aggregation, 3a$\to$5      & \multicolumn{2}{c}{1.88$\times$} & \multicolumn{2}{c}{2.33$\times$} & \multicolumn{2}{c}{2.51$\times$} \\
Capacities, 5$\to$7        & \multicolumn{2}{c}{1.07$\times$} & \multicolumn{2}{c}{1.31$\times$} & \multicolumn{2}{c}{1.55$\times$} \\
vs.\ \naive, 2$'\to$7      & \multicolumn{2}{c}{2.37$\times$} & \multicolumn{2}{c}{3.55$\times$} & \multicolumn{2}{c}{4.56$\times$} \\
vs.\ production, 1$\to$7   & \multicolumn{2}{c}{\textbf{4.74$\times$}} & \multicolumn{2}{c}{\textbf{3.65$\times$}} & \multicolumn{2}{c}{\textbf{4.68$\times$}} \\
\bottomrule
\end{tabular}
\end{footnotesize}
\end{center}
\vskip -0.15in
\end{table}

Table~\ref{tab:ladder} gives the results (Figure~\ref{fig:ladder} in Appendix~\ref{app:ladder} plots them). \sys prefills
128-, 256- and 512-token chunks in 105.3, 141.5 and 225.3\,ms, 4.74$\times$, 3.65$\times$ and 4.68$\times$ faster than the production
compilation and 2.37$\times$, 3.55$\times$ and 4.56$\times$ faster than naive \naive. The attention rewrite (2$\to$2$'$)
saves 17, 67 and 130\,ms in every head-parallel program alike; it is not an MoE gain. Without it, naive expert parallelism
is 10\% slower than production at $T=256$ and 512; with it, 2.5--2.7\% faster, and 2.00$\times$ faster at $T=128$, where
production takes about as long as for 256 tokens. The routing scans (3a)
give 1.17--1.18$\times$. MoE output aggregation, steps 3b--5, is the largest factor at 1.88$\times$, 2.33$\times$ and
2.51$\times$, and the ranked capacities add 1.07--1.55$\times$. Only with both does a longer chunk pay: per token, \sys costs
0.82, 0.55 and 0.44\,ms at $T=128$, 256 and 512, while naive \naive stays at 1.95--2.01\,ms.

\subsection{Capacities and Aggregation Together}
\label{sec:eval-abl}

\textbf{The interaction.} Figure~\ref{fig:interaction} isolates it: all four programs rank experts at run time, gather
weights from the same bank and share the final sum and the tile size. The repaired dense combine keeps one zero-filled
accumulator per tier width, never reads its zeros back, lets a width group's later tiers read, add and scatter back, and
sums lanes in 16 token tiles. With it, tiers change latency by $+36.9\%$, $+3.9\%$ and $-3.6\%$; with token-owned
aggregation, by $-8.8\%$, $-13.3\%$ and $-26.9\%$. Token-owned aggregation is 1.27--1.36$\times$ faster than the dense
combine at capacity $T$, and 1.63--1.90$\times$ with tiers.

\begin{table}[t]
\caption{Capacity attribution (ms per chunk, each program at its best tile size), and expert rows provisioned per SoC.}
\label{tab:attr}
\vskip 0.1in
\begin{center}
\begin{footnotesize}
\setlength{\tabcolsep}{3pt}
\begin{tabular}{@{}lrrr@{}}
\toprule
 & $T=128$ & $T=256$ & $T=512$ \\
\midrule
A static, fixed order, $T$       & 113.8 & 168.7 & 288.0 \\
B gathered, fixed order, $T$     & 109.1 & 160.0 & 302.3 \\
C gathered, ranked, $T$          & 115.5 & 163.1 & 308.1 \\
D ranked, two tiers              & 106.3 & \textbf{138.0} & 269.6 \\
E \sys, three tiers              & \textbf{105.3} & 141.5 & \textbf{225.3} \\
\midrule
C/E, capacities                    & 1.10$\times$ & 1.15$\times$ & 1.37$\times$ \\
A/E, vs.\ best static            & 1.08$\times$ & 1.19$\times$ & 1.28$\times$ \\
Rows, A--C/E                     & 4096/1664 & 8192/3072 & 16384/6144 \\
\bottomrule
\end{tabular}
\end{footnotesize}
\end{center}
\vskip -0.15in
\end{table}

\textbf{Attributing the capacities.} Steps 5--7 change weight access, the assignment of experts to lanes and the
capacities at once. Table~\ref{tab:attr} separates them with five programs that share the aggregation: (A) a fixed expert
order at capacity $T$ with static weights; (B) the same order through the bank gathers; (C) the run-time rank at capacity
$T$; (D) two tiers; (E) \sys. All give bit-identical logits. With ranking and indirection matched, \sys's capacities are
1.10, 1.15 and 1.37$\times$ faster (C$\to$E), and \sys is 1.08--1.28$\times$ faster than the best static program. Each
program runs at its best tile size, the granularity at which the compiler skips padded rows (\S\ref{sec:skip},
Appendix~\ref{app:attrprof}): with 512\,KiB tiles the static program computes a quarter of its rows, while the gathered
programs spill activations at that size and use 1024\,KiB tiles at $T=512$.

\subsection{What Makes the Aggregation Efficient}
\label{sec:eval-own}

\begin{table}[t]
\caption{Aggregation formulations at capacity $T$ with static weights and the same routing: full-model latency relative to
\sys's aggregation (ladder step 5), cross-SoC MiB per MoE layer at $T=512$ (layer 1 of instrumented two-layer cuts), and
whether the logits equal step 5's or differ in FP16 summation order.}
\label{tab:own}
\vskip 0.1in
\begin{center}
\begin{footnotesize}
\setlength{\tabcolsep}{2.5pt}
\begin{tabular}{@{}lrrrrl@{}}
\toprule
Formulation & 128 & 256 & 512 & MiB & Logits \\
\midrule
Step 5: SoC batch axis        & 1.00 & 1.00 & 1.00 & 14 & -- \\
Local top-$k$, SoC slices     & 3.44 & 4.04 & 4.14 & 1539 & equal \\
Local top-$k$, global rows    & 4.37 & 5.14 & 5.30 & 1539 & equal \\
Top-$k$ rows from all SoCs    & 4.05 & 4.70 & 4.92 & 1538 & order \\
Lane-indexed, lane axis       & 1.12 & 1.11 & 1.15 & 12 & order \\
Lane-indexed, SoC slices      & 7.02 & 8.22 & 4.31 & 577 & order \\
Repaired dense combine        & 1.25 & 1.57 & 1.64 & 48 & order \\
\bottomrule
\end{tabular}
\end{footnotesize}
\end{center}
\vskip -0.15in
\end{table}

Table~\ref{tab:own} compares \sys's aggregation, ladder step 5, with six formulations at the same capacity and routing.
The two local top-$k$ variants compute exactly its sum in its order and differ only in how the graph expresses device
ownership; both make the compiler ship expert outputs between SoCs, as does the GPU-style gather of each token's $k$ rows
from all SoCs (\S\ref{sec:mot-own}). The repaired dense combine keeps locality, but its accumulator grows with $T$:
replacing the representation (step 4) gives 1.23--1.63$\times$ on top of the 1.49--1.54$\times$ of the repairs (3b).

\textbf{Local reduction.} The lane-indexed formulation keeps device locality, with slightly less cross-SoC traffic than
step 5 (12 against 14\,MiB), and is still 11--15\% slower. It sums each token tile over a SoC's 16 lanes, which sit on 16
cores: in the profile this sum moves 242\,MiB between cores per layer and keeps them busy for 47\,core-ms, against 1.9\,MiB
and 2.4\,core-ms for step 5's sum over each token's $k$ rows on one core, and its 16 tiles run largely one after another.
The window from the last expert GEMM to the end of the final sum grows from 1.2--1.7 to 3.1--3.7\,ms per SoC
(Appendix~\ref{app:phases}); the uninstrumented full models differ by 1.1\,ms per layer, so these windows locate the
difference rather than decompose it along a dependency chain.

\subsection{Tier Structure}
\label{sec:eval-pad}

Each tier is a stage with a cost of its own. In a full-model profile of \sys at $T=512$ (before the attention fix), the
8$\times$128 and 16$\times$64 tiers span 0.59 and 0.52\,ms per layer while their GEMMs keep the busiest core busy for
0.08--0.09\,ms: they wait on their experts' weights, and each tier starts after a gap of 0.28--0.48\,ms. The first tier
splits because it is bound by the SoC's memory traffic: 16 lanes at capacity 512 take 2.2\,ms, eight at 512 and eight at 128
take 0.8--1.0 and 0.65--0.72\,ms (Appendix~\ref{app:stages}). Three tiers save 28--33\% of the rows of two, and their worth
depends on the chunk: 1.20$\times$ at $T=512$, 1.01$\times$ at 128 and 0.975$\times$ at 256 (Table~\ref{tab:attr}), where a
session before the attention fix had found them 3.6\% faster; four tiers lost to three on two-layer cuts
(Appendix~\ref{app:neg}). Rows do not predict the best structure, and the deployed one is not the best at every chunk length;
selecting it from measured stage costs remains open.

% tab:transfer belongs to sections/transfer.tex, but a table* is set no earlier than the next page, so it is defined
% here to land on the page where that section starts.
\begin{table*}[t]
\caption{What transfers: for each part of \sys, the reasoning we expect to hold on other statically partitioned targets,
the conditions under which it pays, and the Cloud AI 100 compiler behavior that shaped our realization of it.}
\label{tab:transfer}
\vskip 0.1in
\begin{center}
\begin{footnotesize}
\setlength{\tabcolsep}{4pt}
\begin{tabular}{@{}>{\raggedright\arraybackslash}p{0.12\textwidth}>{\raggedright\arraybackslash}p{0.34\textwidth}>{\raggedright\arraybackslash}p{0.22\textwidth}>{\raggedright\arraybackslash}p{0.26\textwidth}@{}}
\toprule
Part & Transferable reasoning & Pays when & Shaped here by the compiler \\
\midrule
Rank-bound capacities & a lane serves a load rank within a device's experts; the $r$-th largest of $n$ assignments is at
most $n/r$; per-rank maxima were stable across workloads, identities were not & expert work is fixed before routing;
padded rows cost time; many experts per device & no control flow, so a sort and gathers; one gathered constant, replicated;
tile size per chunk length \\
Tier structure & first tier at $T$, lower tiers at $h$ times the calibrated maximum of their first rank; rows traded against
measured per-stage cost; overflows show in the counts and a rerun recovers them exactly & a stage's fixed cost is small
against the rows it removes & widths must divide 16 and stages run in sequence \\
Ownership-preserving aggregation & designed with the capacities; device ownership stated in a form the compiler keeps;
each token's rows summed on the core that owns it; one elementwise cross-device sum; cost grows with $kT$ per tier & the
baseline combine is dense in lanes; capacities change the outputs' shapes & index values and slices lost locality, the
SoC batch axis kept it; gather without zero fill \\
\bottomrule
\end{tabular}
\end{footnotesize}
\end{center}
\vskip -0.15in
\end{table*}

\subsection{Fixed-Size Blocks Instead of Ranked Lanes}
\label{sec:eval-blocks}

Block scheduling packs assignments into blocks of $B$ rows, one expert per block, with $\lceil kT/B\rceil+E-1$ blocks to
drop nothing \citep{awsneuronmoe}. Built from \sys's own components and measured on two-layer cuts before the attention fix
(Appendix~\ref{app:blocks}), the best calibrated block program is 1.26, 1.39 and
1.69$\times$ slower than \sys at $T=512$, 256 and 128, the best dropless one 1.51--1.71$\times$, and even a
trace-specialized diagnostic 1.12--1.31$\times$. Blocks lose on weight traffic: every block slot streams its expert's MXFP6
weights, used or not, 337.5\,MiB per SoC and layer for 64-row dropless blocks against 95--106\,MiB for \sys, and one weight
load cannot serve consecutive blocks of an expert in this compiler.

\subsection{Beyond the First Chunk, and Overflow}
\label{sec:eval-chunks}

The programs keep the cache for $2T$ positions. A second chunk attends to the cache the first wrote: \sys and the static
capacity-$T$ program give bit-identical logits for both chunks at all three lengths, and the second chunk costs what the
first does. To test the recovery of \S\ref{sec:overflow}, we forced an undersized capacity profile on the second of four
128-token chunks: the counts flagged 45 (layer, lane) pairs, the rejected logits were 0.04 relative $L_2$ off, and after
the rerun all four chunks were bit-identical to runs without overflow, as after an overflow in the first chunk. Capacity
inputs and extra profiles make the tiered program 18\% slower, so a resident capacity-$T$ program may be the cheaper
fallback, at the cost of re-prefilling earlier chunks, which we did not time. These tests predate the attention fix.

\subsection{Do Calibrated Capacities Generalize?}
\label{sec:eval-gen}

We split each of eight workloads into calibration and evaluation documents before cutting chunks, and measured the routing
of every chunk on the device. \sys's capacities overflowed in none of the 805, 618 and 449 evaluation chunks at $T=128$, 256
and 512 and in no workload (30 chunks hold documents the deployed capacities were calibrated on), nor did
leave-one-workload-out calibration; without headroom, the rule gives a 112-row middle tier at $T=512$, which overflows in 2
of 449 chunks. The same tiers bound to a fixed expert order need capacity $T$ everywhere; calibrated per workload, they need
1.4--2.7$\times$ the rows and still overflow in up to 30\% of chunks, while ranks 1, 9 and 17 never exceeded 509, 119 and 52
tokens. The bound $n_{(r)}\le N/r$ holds with room (at most 0.97 of it).

\subsection{Accuracy}
\label{sec:eval-acc}

Every program stays at the error level of the MXFP6 weights: the relative $L_2$ error of the last-position logits against
FP32 is 0.188, 0.091 and 0.096 for \sys at $T=128$, 256 and 512, against 0.188, 0.089 and 0.095 for production
(Appendix~\ref{app:acc}). The programs differ from production only in FP16 summation order, in attention's group sum and in
the order of a token's expert rows; A--E are bit-identical to each other. The next token of \sys and of production
matches FP32 except in the 256-token chunk, whose two leading FP32 logits are 0.58 apart. No lane exceeded its capacity.

\subsection{A Fused Kernel in the Same Layout}
\label{sec:eval-kernel}

TokenCombine, a C++ custom operator with HVX vector intrinsics, replaces step 5's gather, mask and sum over the same
$[4,16C,d]$ layout (Appendix~\ref{app:custom}); its tiles follow the SoC axis over all 64 cores, and cross-SoC traffic stays
at step 5's 14\,MiB. At $T=512$ it runs at 0.940$\times$ of step 5
in the fit (323.3--324.7 against 341.9--343.8\,ms in one alternating session). Its logits differ deterministically from step
5's by 1.7\% relative $L_2$, because it sums in FP32 in another order, and it takes 650--656 minutes to compile. The layout
thus also supports a fused implementation, with a modest gain at the one chunk length we tested.

\subsection{Limitations}
\label{sec:limits}

The lower tiers are calibrated, not dropless: no evaluation chunk overflowed, but the deployed rate depends on model and
workload. We evaluate one model, at batch 1 and in prefill only (our rewrites fix $T$), and blocks under this compiler
only. The tier structure was chosen by comparison, not by a validated method,
and the ownership results concern this compiler's partitioner. The SoCs share a 150\,W limit under which most padded
programs throttle at times (\sys never does), so we compare programs through the anchored fit rather than across sessions.
